
\section{Conclusions and Future Work}
\label{sec:5}

Vehicular Ad Hoc Networks (VANETs) are an
important component of Intelligent
Transportation Systems (ITS), enabling
communication among vehicles, roadside
infrastructure, and other transportation
entities. These communication capabilities
support applications related to road safety,
traffic management, cooperative driving,
and transportation efficiency.

This paper presented a literature-based
review of VANET architectures, communication
models, routing protocols, simulation
frameworks, performance evaluation methods,
and emerging V2X technologies. The review
examined Vehicle-to-Vehicle (V2V),
Vehicle-to-Infrastructure (V2I),
Vehicle-to-Pedestrian (V2P), and
Vehicle-to-Network (V2N) communication
as important components of the broader
Vehicle-to-Everything (V2X) concept.

The analysis of routing protocols highlighted
the trade-offs among proactive, reactive,
and geographic routing approaches.
Proactive protocols maintain routing
information in advance but may introduce
continuous control overhead. Reactive
protocols discover routes when needed,
which can reduce unnecessary route
maintenance but may introduce route
discovery delays. Geographic routing
uses location information to support
forwarding decisions, although its
effectiveness depends on positioning
accuracy, connectivity, and road
topology.

The comparative analysis of published
studies further demonstrated that
routing performance cannot be assessed
independently of the operating
environment. Vehicle density, mobility
patterns, road layouts, communication
conditions, and application requirements
can influence the relative suitability
of different approaches. Therefore,
no single routing protocol should be
considered universally optimal for
all VANET scenarios.

This paper also reviewed simulation
and evaluation tools, including SUMO,
OMNeT++, Veins, and NS-3. Performance
metrics such as packet delivery ratio,
end-to-end delay, throughput, packet
loss ratio, and routing overhead were
discussed as important indicators of
network behavior. The review emphasized
that meaningful comparisons require
consistent simulation assumptions,
realistic mobility models, and clearly
defined evaluation metrics.

To complement the literature review,
a simplified Python-based illustrative
experiment was implemented to examine
packet delivery ratio (PDR) and average
communication delay under different
vehicle counts.

The experiment evaluated five scenarios
with 10, 20, 30, 40, and 50 vehicles
using a simplified probabilistic
communication model.

The generated PDR values ranged from
21.80\% to 37.20\%, while the estimated
average communication delay increased
from 15.33 ms for 10 vehicles to
35.97 ms for 50 vehicles.

The results were presented using a
numerical table and two graphical
visualizations, illustrating how
changes in vehicle count can influence
the performance indicators produced
by the selected model.

However, these results are dependent
on the assumptions and mathematical
rules used in the experiment.
The implementation does not include
realistic vehicle mobility, detailed
wireless channel behavior, or
established VANET routing protocols.
Therefore, the results should be
interpreted as an educational
illustration rather than validated
measurements of real-world
vehicular network performance.

Furthermore, the evolution of vehicular
communication technologies was examined
through IEEE 802.11p-based systems,
LTE-V2X, IEEE 802.11bd, and 5G NR-V2X.
These technologies offer different
communication mechanisms and resource
management approaches. Their suitability
depends on deployment conditions,
application requirements, infrastructure
availability, and interoperability
considerations rather than technological
generation alone.

Despite continued development, VANETs
face several unresolved challenges.
High mobility, unstable connectivity,
wireless channel congestion, security
attacks, privacy concerns, and
heterogeneous technology deployment
remain important obstacles to
reliable operation.

Future research should investigate
mobility-aware routing, congestion
management, secure and privacy-preserving
communication, and interoperability
between different V2X technologies.
Edge computing and artificial
intelligence may also provide
opportunities for adaptive resource
management and network optimization.
However, such approaches require
careful evaluation of computational
overhead, generalization, security,
and practical deployment constraints.

An additional research priority is
the development of reproducible
evaluation methodologies that combine
realistic traffic mobility,
wireless communication models,
and clearly documented experimental
configurations. Such methodologies
can improve the reliability and
comparability of findings reported
across different studies.

Future work may also extend the
illustrative Python experiment by
integrating realistic vehicle
mobility using SUMO and detailed
wireless network simulation using
NS-3 or OMNeT++.

Such an extension could support
more realistic investigations of
packet delivery, communication
delay, channel congestion, and
routing protocol performance
under different traffic conditions.

Overall, this paper combines a
structured literature-based review
with a reproducible illustrative
Python experiment. Its contribution
lies in organizing established
research, comparing important
technical approaches, identifying
open challenges, and demonstrating
basic performance evaluation concepts.

VANETs provide a foundation
for increasingly connected
transportation systems. Their
successful development will depend
not only on advances in wireless
communication, but also on the
integration of dependable routing,
security, privacy protection,
realistic evaluation, and
practical deployment strategies.
